\subsection{Hydrodynamic Plausibility of the Generated Curves}
\label{subsec:agit-plausibility}

The curves generated by the apparatus are compared with three references: a CFD simulation of each impeller, the Power-number curves measured by students on an Armfield stirred-vessel unit for three other impellers (a stock paddle, a custom paddle, and custom blades), and the literature curve tabulated by Brown et al.\ (curve~10 of the reference, here the ``literature'' curve). % TODO: add the BibTeX key of Brown et al. and of the Armfield student report, and give the Armfield unit and vessel geometry.
All Reynolds and Power numbers use Equations~\ref{eq:npo-torque} with $\rho = 996.3$~kg\,m$^{-3}$, $\mu = 0.854$~mPa\,s and the diameters of the dashboard, $D = 0.0665$~m for the 96~mm impeller and $0.1207$~m for the 174~mm impeller. The CFD is available as the torque at each simulated speed (55 speeds for each impeller, 0.01 to 11\,600~rpm for the 96~mm and 0.004 to 3500~rpm for the 174~mm impeller) and is converted with the same equations. The tank was baffled in every experimental run and in the simulation. The liquid level was 200~mm, with the bottom of the impeller 50~mm above the tank floor for the 96~mm impeller and 90~mm for the 174~mm impeller.

\paragraph{Data and method.}
Each impeller was run several times through the dashboard sequence (28 setpoints for the 96~mm impeller, 23 for the 174~mm impeller, 15~s per setpoint). One run per impeller is used for the comparison, the best one, chosen by a rule fixed before the comparison: the run must be complete; the torque must not fall while the speed rises at $N_{Re} \geq 5\times10^{4}$ (a collapse of the torque, as caused by power limiting or by air drawn into the liquid); and among the runs that pass, the one with the smallest coefficient of variation of $N_{Po}$ over the plateau $N_{Re} \geq 5\times10^{4}$ is taken. The best run is that of 6 November 19:17 for the 96~mm impeller (variation 1.4~\%) and that of 6 November 21:17 for the 174~mm impeller (3.7~\%). Figure~\ref{fig:runs} shows it against the other complete runs. One run, of 6 November 19:31, has a flatter plateau (3.5~\%) but is excluded because its torque falls at the highest speeds. The Power number at the plateau of the other complete runs ranges from 1.69 to 2.09 for the 96~mm and from 2.12 to 2.53 for the 174~mm impeller, so the best run is at the upper end of the 174~mm runs, and the comparison below is for that run and not for the typical run. % TODO: confirm the cause of the runs that rise slowly (96 mm: 5 Nov 22:48, 6 Nov 19:09 and 22:09) and of the torque collapse (174 mm, 6 Nov 19:45 and 19:53) from the run notes, and the torque zero offset of each run.
Torque below 0.002~N\,m, about two counts of the sensor, is treated as noise and omitted. The CFD torque is interpolated linearly in the logarithms of speed and torque to each experimental speed. The error measure is the root-mean-square error (RMSE)
\begin{equation}
    \mathrm{RMSE} = \sqrt{\frac{1}{n}\sum_{i=1}^{n}\left(y_{\mathrm{ref},i} - y_{\mathrm{data},i}\right)^{2}}, \qquad \mathrm{NRMSE} = \frac{\mathrm{RMSE}}{\overline{y}_{\mathrm{data}}}
    \label{eq:rmse}
\end{equation}
of the torque (N\,m) or of $N_{Po}$, with $y_{\mathrm{ref}}$ the reference (CFD or literature) interpolated to the Reynolds number of each point of the data, over the points in which both exist. Because the CFD is a trusted reference only in the turbulent range, the main comparison uses $N_{Re} \geq 4\times10^{4}$; the comparison over every point above the sensor floor is also given. The error bars in the figures are the standard deviation of the torque within each averaging window of the run.

\begin{figure}[htbp]
    \centering
    \includegraphics[width=\textwidth]{fig_runs}
    \caption[All complete runs and the best run]{Power number against Reynolds number for all complete runs of the 96~mm (left) and 174~mm (right) impeller. The best run (coloured) is chosen by the rule in the text; the others are grey.}
    \label{fig:runs}
\end{figure}

\paragraph{Torque and Power number against the CFD.}
Figures~\ref{fig:cfd-96} and~\ref{fig:cfd-174} compare the best run of each impeller with the CFD, point by point, and Table~\ref{tab:rmse-cfd-mixr} gives the error. In the range $N_{Re} \geq 4\times10^{4}$ the CFD underpredicts the torque of the 96~mm impeller by 27~\% on average ($\tau_{\mathrm{CFD}}/\tau_{\mathrm{exp}} = 0.73$, RMSE 0.038~N\,m) and that of the 174~mm impeller by 18~\% (0.82, RMSE 0.048~N\,m). The Power number of the experiment is 2.02 for the 96~mm and 2.56 for the 174~mm impeller, against 1.46 and 2.08 for the CFD, an RMSE of 0.59 (NRMSE 29~\%) and 0.50 (20~\%). Below $N_{Re} = 4\times10^{4}$ the experimental points scatter widely and fall below the CFD, as the torque there is a few times the resolution of the sensor (panel~c of both figures). The ratio is nearly constant above $4\times10^{4}$ for both impellers, so the CFD and the experiment have the same dependence on speed, and the disagreement is in the level of the torque, $-18$ to $-27$~\%. The CFD is baffled like the experiment, so a difference in the baffling does not explain the gap. The conditions of the run, the zero of the torque sensor, and the geometry of the simulated impeller, which are not tested here, remain as causes.

\begin{table}[htbp]
    \centering
    \caption[CFD against the best run: torque and Power number]{Error of the CFD against the best run of each impeller in torque and in Power number. Range: ``Re $\geq 4\times10^{4}$'' (Reynolds number) or ``all above sensor floor''. $n$ is the number of setpoints, $\tau_{\mathrm{CFD}}/\tau_{\mathrm{exp}}$ the mean ratio, and the overlined quantities the means over the points. RMSE and NRMSE are defined in Equation~\ref{eq:rmse}.}
    \label{tab:rmse-cfd-mixr}
    \small
    \resizebox{\textwidth}{!}{\input{Chapter5/Figs/best_cfd_vs_mixr}}
\end{table}

\begin{figure}[htbp]
    \centering
    \includegraphics[width=\textwidth]{fig_cfd_96mm}
    \caption[96 mm impeller against the CFD]{96~mm impeller, best run (filled markers, error bars: standard deviation of the torque within each averaging window) against the CFD (open circles, one per simulated speed). (a) Power number against Reynolds number, with the literature value (dashed). (b) Torque against speed. (c) Ratio of the CFD torque to the measured torque; open markers: $N_{Re} < 4\times10^{4}$; the shaded band is $\pm10$~\%.}
    \label{fig:cfd-96}
\end{figure}

\begin{figure}[htbp]
    \centering
    \includegraphics[width=\textwidth]{fig_cfd_174mm}
    \caption[174 mm impeller against the CFD]{As Figure~\ref{fig:cfd-96} for the 174~mm impeller.}
    \label{fig:cfd-174}
\end{figure}

\paragraph{CFD against the Armfield student data.}
The students' curves are for different impellers, so the comparison measures how far their curves lie from the CFD curves of this work, and not an error of the CFD. Figure~\ref{fig:armfield} shows them and Table~\ref{tab:rmse-cfd-arm} gives the RMSE over the points with $N_{Re} \geq 4000$, the lowest Reynolds number at which the CFD is used (the first point of the stock and custom paddle, at 3912, is therefore omitted). The stock paddle has $N_{Po}$ of 7.8 to 12.1 over $N_{Re} = 1.5\times10^{4}$ to $3.4\times10^{4}$, about six times the CFD, with an RMSE of 8.2 against the 96~mm and 7.7 against the 174~mm CFD. The custom paddle falls from 1.6 at $1.8\times10^{4}$ to 0.4 at $8\times10^{4}$ and the custom blades from 8.0 at $7.9\times10^{3}$ to 0.2 at $8\times10^{4}$. Their RMSE is 0.84 and 1.42 for the paddle, and 2.2 and 2.4 for the blades, against the 96 and 174~mm CFD, which is 116 to 195~\% of their mean $N_{Po}$. The RMSE of the blades is dominated by their lowest-$N_{Re}$ points: it falls to 1.1 against the 96~mm CFD when only the points above $2\times10^{4}$ are used. The students' $N_{Po}$ keeps falling with $N_{Re}$ up to $8\times10^{4}$, whereas both CFD curves and the literature curve are flat above about $10^{4}$. A Power number that falls by more than a decade over a decade of Reynolds number is not the behaviour of fully turbulent flow, and the more likely reason is a torque reading that is small against the zero offset of the measurement at the low speeds of that unit, the same cause that the MIXR-1 data show at low speed (Section~\ref{subsec:agit-torque-drift}). That explanation is not tested here.

\begin{figure}[htbp]
    \centering
    \includegraphics[width=\textwidth]{fig_armfield}
    \caption[Armfield student data against the CFD]{Armfield student data (stock paddle, custom paddle, custom blades), the CFD of the 96~mm (left) and 174~mm (right) impeller, and the literature curve, on logarithmic axes.}
    \label{fig:armfield}
\end{figure}

\begin{table}[htbp]
    \centering
    \caption[CFD against the Armfield student data]{Error of the CFD (the 96 and 174~mm impellers of this work) against the three Armfield student curves, for the points with $N_{Re} \geq 4000$. $n$ is the number of points, the overlined quantities are means over the points, and the last column is the RMSE of $\log_{10}N_{Po}$, a measure that does not weight the large low-$N_{Re}$ values.}
    \label{tab:rmse-cfd-arm}
    \small
    \resizebox{\textwidth}{!}{\input{Chapter5/Figs/rmse_cfd_vs_armfield}}
\end{table}

\paragraph{Comparison with the literature curve.}
The literature curve is flat at $N_{Po} \approx 1.7$ to 1.8 over $N_{Re} = 4\times10^{3}$ to $10^{6}$. Figure~\ref{fig:literature} places the best runs and the CFD on it, and Table~\ref{tab:rmse-lit} gives the error of every series over the Reynolds numbers where both exist. The best 96~mm run is 14~\% above the literature (RMSE 0.31) and the best 174~mm run is 45~\% above it (RMSE 0.81); the CFD is 18~\% below it for the 96~mm and 18~\% above it for the 174~mm impeller (RMSE 0.32 and 0.31). The Armfield curves are far from it: the stock paddle by a factor of 5.4 (RMSE 8.0), the custom paddle 59~\% below it (RMSE 1.1) and the custom blades 19~\% below it on average (RMSE 2.4; the mean bias understates the error because the large positive errors at low $N_{Re}$ cancel the negative errors at high $N_{Re}$). The CFD and the 96~mm data therefore agree with the literature to within 14 to 18~\%, the 174~mm data are further from it, and the student curves do not agree with any of them. The geometry of the MIXR-1 impellers is not that of the standard configuration of the literature curve, so the offset is not attributed to the measurement. % TODO: state the geometry of the reference configuration (impeller type, D/T, C/T, number of baffles) and in which ratios the 96 and 174~mm impellers depart from it.

\begin{figure}[htbp]
    \centering
    \includegraphics[width=0.7\textwidth]{fig_literature}
    \caption[Power number against the literature curve]{Best run of each impeller (markers) and CFD (open circles) against the literature curve (dashed).}
    \label{fig:literature}
\end{figure}

\begin{table}[htbp]
    \centering
    \caption[Every series against the literature curve]{Error of each series against the literature curve (Brown et al., curve~10), interpolated in the logarithms to the Reynolds number of each point. MIXR-1 and CFD: $N_{Re}\geq4\times10^{4}$ and, for the CFD, up to the highest experimental value; Armfield: $N_{Re}\geq4000$. Bias is the mean of $(N_{Po} - N_{Po,\mathrm{lit}})/N_{Po,\mathrm{lit}}$.}
    \label{tab:rmse-lit}
    \small
    \resizebox{\textwidth}{!}{\input{Chapter5/Figs/best_vs_literature}}
\end{table}

\paragraph{Limits of the comparison.}
Four limits apply. The comparison uses one run per impeller, the best by the rule above, and the other runs differ from it by up to 20~\% in $N_{Po}$ at the plateau; the conditions of the run (the zero of the torque sensor and the state of the liquid) are not recorded. The liquid level is the same for the two impellers but the clearance is not (50~mm and 90~mm). The Armfield and literature data were transcribed from a table, so their Reynolds and Power numbers carry the rounding of that table. And the Reynolds number range of the apparatus ends at $1.2\times10^{5}$, so the comparison with the literature curve says nothing about the laminar and transitional range, where the stock paddle and the CFD show the largest changes.

